The reviewed literature spans fault taxonomies, dependency and propagation modelling, cross-layer telemetry, resilience patterns, and runtime \gls{fdir} strategies across \gls{soa}, microservices, cloud, IoT, aerospace, automotive, and embedded systems. These research areas correspond directly to the guiding questions introduced in Chapter~\ref{chap:methodology}, which investigate existing fault-management solutions, fault detection mechanisms, dependency modelling approaches, and recovery strategies in service-oriented systems. Although the individual contributions differ in scope and domain, a coherent set of conceptual pillars emerges that together shape the current understanding of fault management in distributed, service oriented, and resource-constrained environments.

The literature is assessed using four criteria motivated by industrial router constraints: (i) whether an approach supports runtime fault episode tracking (distinguishing new fault occurrences from ongoing violations), (ii) whether recovery decisions are dependency-aware (root cause vs.\ downstream symptoms), (iii) whether recovery execution is coordinated and gated (avoiding repeated actions for the same episode), and (iv) whether the approach offers deterministic and bounded overhead suitable for embedded platforms. These criteria are used throughout Sections 4.1 to 4.6 to identify which capabilities are available in existing work and which gaps remain.

\section{Fault Models and Recovery Oriented Classification}\label{sec:lit-fault-models}

A clear pattern across the taxonomy focused works is the shift from detection centric classifications to recovery oriented, multilevel models. Recovery coupled taxonomies in \cite{ref1,ref2} enrich fault classification with severity, locality, priority, and recoverability attributes, providing actionable semantics for downstream recovery logic. Broader surveys such as \cite{ref3,ref4} extend this with sensing, actuator, network, semantic, and composition level fault categories, reflecting the diversity of fault origins in adaptive, \gls{soa}, and cloud systems. The SOC specific analyses in \cite{ref5,ref6} expose additional classes relevant to service contracts, control flow, and composition semantics.

While recovery-oriented taxonomies provide actionable attributes (e.g., severity, locality, recoverability), they typically remain classification frameworks and do not define how repeated observations over time should be grouped into a single operational incident. As a result, they do not provide explicit runtime support for fault episode tracking or for gating recovery actions across repeated violations in long-running, resource-constrained router systems.

\section{Dependency Modelling and Error Propagation}\label{sec:lit-dependency-modelling}

A second recurrent theme concerns explicit modelling of how erroneous states traverse software and system boundaries. Quantitative propagation concepts such as error propagation probabilities in \cite{ref7}, error permeability and exposure from \cite{ref8}, and propagation aware reliability formulations in \cite{ref9} demonstrate that system reliability fundamentally depends on how faults propagate not just on the failure probability of isolated components. Formal dependency structures including CPN derived control/data dependencies in \cite{ref10}, service/task graphs in \cite{ref11}, and Bayesian networks for cross layer causal propagation in \cite{ref12} further reinforce the importance of explicit dependency modelling.

These works converge on the insight that fault management requires a representation of structural and behavioural dependencies, augmented with propagation semantics. However, the available models are either design time analytical tools or runtime systems that exceed the computational envelope of router class devices. None articulate a a router-specific dependency model with bounded runtime and memory costs that captures service to service, service to network, and network to service propagation paths with bounded and predictable runtime overhead. Although these approaches establish the importance of dependency-aware fault reasoning, they primarily address diagnosis or reliability analysis in isolation and do not integrate dependency traversal with coordinated, episode-aware recovery execution under deterministic runtime constraints. Existing dependency and propagation models either target offline reliability analysis or use probabilistic/formal techniques that are difficult to apply with bounded runtime overhead on router-class devices. Consequently, they do not provide a practical runtime mechanism that couples dependency traversal with episode-aware recovery coordination (i.e., selecting a root cause while suppressing downstream symptoms during the same fault occurrence).

\section{Telemetry, Observability, and Cross Layer Fault Information}\label{sec:lit-telemetry}

The literature consistently demonstrates that effective fault localisation and recovery depend on comprehensive, correlated telemetry across layers. Cross layer telemetry integration in \cite{ref13} shows how embedding path level network metrics into application level traces accelerates causal diagnosis. Bandwidth aware telemetry prioritisation and in agent degradation handling in \cite{ref14} illustrate the need for resilient telemetry channels even under control plane stress. Event based cross layer signalling in \cite{ref15} highlights the importance of a unified fault event communication substrate. Studies on failure amplification in constrained networks, such as \cite{ref16}, underscore that faults in communication substrates directly influence system level behaviour.

Despite these advances, none of the reviewed approaches deliver a telemetry model specifically designed for industrial routers and their \gls{soa} resident services. Router constraints limited \gls{cpu}/memory, bounded control plane bandwidth, and the need for deterministic timing are absent from current telemetry frameworks, which limits their applicability in embedded, safety critical networking platforms. As a result, existing telemetry frameworks lack explicit mechanisms for stateful fault episode tracking and controlled fault correlation suitable for industrial routers, where telemetry processing must remain bounded, deterministic, and tightly coupled to recovery logic. Although cross-layer observability improves diagnosis, most telemetry frameworks focus on data collection and correlation rather than on maintaining fault episode state that governs recovery execution. For industrial routers, the missing capability is a telemetry model that supports deterministic, bounded processing while explicitly enabling episode-level tracking and recovery gating at runtime.

\section{Runtime \gls{fdir}, Self Healing, and Deterministic Recovery}\label{sec:lit-fdir-recovery}

A significant body of work focuses on runtime \gls{fdir} and self healing in microservices, fog/edge, IoT, aerospace, and automotive systems. Pattern driven resilience architectures such as \cite{ref17} and orchestrator based recovery models in \cite{ref18,ref19} emphasise structured, repeatable, and explainable recovery paths. Reactive and locality aware edge/fog strategies, exemplified by \cite{ref20}, demonstrate efficient task and node level recovery in distributed environments. Multiagent fault recovery designs (\cite{ref21,ref22}) highlight decentralised coordination of diagnosis and healing. Fail operational embedded strategies, including \cite{ref23,ref24,ref25}, provide evidence for deterministic failover under strict timing constraints. \gls{tsn} supported resilient \gls{soa} prototypes such as \cite{ref26} integrate service orchestration, redundancy, and monitoring across time sensitive networks.

These works collectively offer a rich catalogue of recovery strategies, resilience patterns, and orchestration mechanisms. Their primary limitation, however, is the assumption of resource rich environments (cloud, microservices platforms) or domain specific constraints (AUTOSAR, \gls{tsn} demonstrators). While these works demonstrate effective recovery mechanisms and resilience patterns, they do not provide a unified runtime control loop that maintains fault episode state and coordinates recovery that coordinates root cause identification, executes recovery actions once per fault episode, and suppresses dependent symptom faults within the resource and determinism constraints of industrial router platforms. Runtime self-healing frameworks provide rich recovery patterns and orchestration mechanisms, but they rarely make fault episode state a first-class runtime concept that constrains when recovery may execute. As a result, they do not systematically enforce “execute once per fault occurrence” behavior nor coordinate symptom suppression with root-cause recovery under embedded resource constraints.

\section{\gls{ml} Based Troubleshooting Versus Deterministic Methods}\label{sec:lit-ml-vs-deterministic}

\gls{ml} driven fault prediction and causal diagnosis approaches including \cite{ref27,ref28}, and the methods surveyed in \cite{ref6} demonstrate high accuracy in cloud scale service environments, particularly when dense telemetry and stable dependency structures are available. However, these techniques rely on computationally intensive models, large training datasets, and assumptions of nondeterministic latency, which are incompatible with deterministic embedded platforms. Their value for industrial routers lies mainly in structural insights (e.g., propagation direction heuristics, metric correlation patterns), not in direct applicability. Consequently, while \gls{ml}-based methods offer valuable insights for large-scale environments, they do not address the need for deterministic, episode-aware fault management with predictable resource usage and explainable decision paths required in embedded industrial routers. \gls{ml}-based diagnosis can be effective when abundant telemetry and compute resources are available, but it does not meet the transparency and bounded-overhead requirements typical of industrial routers. In particular, these approaches do not directly address deterministic, episode-aware recovery coordination, where recovery actions must be explainable and consistently gated across repeated observations.s

\section{Cross Layer Integration and Service Oriented Diagnostics}\label{sec:lit-cross-layer}

Integration oriented works such as \cite{ref29,ref15}, together with multilayer frameworks like \cite{ref2}, demonstrate how diagnostic and \gls{fdir} components can be modularised as services and coordinated across layers. Time-Sensitive Networking (\gls{tsn}), a set of IEEE 802.1 standards that provide deterministic, bounded-latency communication over Ethernet, is relevant in this context because industrial routers increasingly support \gls{tsn}-capable network planes alongside conventional \gls{ip} forwarding. Ergenc et al.\ \cite{ref26} demonstrate a \gls{tsn}-enabled \gls{soa} prototype with service migration and redundancy, but
their work does not address service-dependency-aware episode management or stateful recovery coordination. They show that fault management benefits from unifying detection, dependency reasoning, and recovery across diverse architectural layers. Yet, no existing work combines these principles into a router specific architecture capable of managing routing plane, \gls{tsn} plane, and \gls{soa} service faults as part of a single cohesive model. Existing integration-oriented approaches thus fall short of delivering a router-specific, stateful \gls{fdir} architecture that unifies cross-layer diagnosis with dependency-coordinated, episode-gated recovery and impact-aware symptom management. Integration-oriented approaches demonstrate that diagnostics and recovery can be composed across layers, but they do not combine cross-layer integration with explicit fault episode tracking and dependency-coordinated recovery gating in a router-specific setting. The remaining gap is therefore a unified runtime model that ties together classification, telemetry, dependency reasoning, and episode-aware recovery coordination.

\section{Summary}\label{sec:lit-summary}

Overall, the reviewed literature provides mature and reusable foundations for fault management in distributed and service-oriented systems:

\begin{itemize}
\item Recovery oriented fault taxonomies with actionable semantics (\cite{ref1,ref2})
\item Quantitative dependency and propagation models (\cite{ref8,ref7,ref9})
\item Formal service dependency structures (\cite{ref11,ref10,ref12})
\item Advanced cross layer telemetry and fault signalling (\cite{ref13,ref14,ref15})
\item Rich resilience and recovery patterns (\cite{ref17,ref18,ref19})
\item Deterministic failover models in embedded \gls{soa} like environments (\cite{ref25,ref23,ref24})
\end{itemize}

Despite this breadth, the literature addresses these aspects largely in isolation or under assumptions that do not hold for long-lived, resource-constrained industrial routers. In particular, existing approaches typically focus on classification, diagnosis, observability, or recovery as separate concerns, and rarely maintain explicit runtime state that distinguishes new fault occurrences from ongoing violations across time.

Across the reviewed work, four recurring limitations emerge:

\begin{enumerate}
\item Fault taxonomies provide expressive classification schemes but do not support runtime differentiation between initial fault occurrence and repeated observations, limiting their ability to govern recovery execution over time.
\item Telemetry and observability frameworks emphasise data collection and correlation but do not explicitly encode fault episode state required for deterministic recovery gating under bounded resource constraints.
\item Dependency and propagation models enable diagnosis and reliability analysis but are primarily used offline or for probabilistic reasoning, rather than for coordinating runtime root-cause recovery and suppressing downstream symptoms.
\item Runtime recovery and self-healing frameworks offer rich resilience patterns but do not consistently enforce episode-aware recovery semantics, such as executing corrective actions at most once per fault occurrence while preventing collateral recovery actions in dependent services.
\end{enumerate}

As a result, existing systems remain vulnerable to repeated recovery execution, cascading alerts and repeated recovery actions, and inconsistent handling of cascading faults in industrial router environments.

This gap motivates the contribution of this thesis: the design of a runtime-stateful, deterministic Fault Detection, Isolation, and Recovery (\gls{fdir}) framework for \gls{soa}-based industrial routers. The framework integrates structured fault characterisation, structured telemetry, explicit dependency modelling, and episode-aware recovery coordination into a single cohesive model. By treating fault episodes as first-class runtime entities, the approach enables root-cause-focused recovery, coordinated suppression of dependent symptoms, and predictable behaviour under embedded system constraints. The following chapter builds on these insights and presents the conceptual architecture of the proposed stateful \gls{fdir} framework for \gls{soa}-based industrial routers.